

\normalsize
\aktspis{}
\akthead
%\begin{dwieszpalty}
\mtyt{Nowa cząstka w~LHC!}


Hadrony to cząstki złożone z~kwarków, które ze względu na swoją budowę dzielą się na dwie główne grupy: bariony oraz mezony. Bariony, takie jak proton ($uud$) i~neutron ($udd$), składają się zawsze z~trzech kwarków. Natomiast mezony, jak na przykład $J/\Psi$ (układ kwarka $c$ i~jego antycząstki $\bar{c}$) albo piony $\pi^+$ ($u\bar{d}$), $\pi^-$~($d\bar{u}$), powstają z~połączenia kwarka i~antykwarka.



Oddziaływania silne, odpowiedzialne za wiązanie kwarków, są na tyle potężne, że kwarki nigdy nie występują samodzielnie. Samotny kwark, posiadający kolorowy ładunek oddziaływań silnych, ulega procesowi \Magenta{\textbf{hadronizacji}}. W~jego trakcie oddziaływanie silne, pośredniczone przez gluony, prowadzi do tworzenia dodatkowych cząstek, które łączą się z~kwarkiem, tworząc hadrony.

Wyjątek stanowi ciężki kwark %$t$
wysoki, nazywany z~angielskiego kwarkiem $t$ (top), co czyni go niewątpliwie najbardziej osobliwym z~kwarków. Tę odkrytą w~1995 roku (w~amerykańskim Tevatronie) cząstkę -- najcięższą ze wszystkich opisywanych przez Model Standardowy -- cechuje czas życia rzędu $10^{-25}$ sekundy. Ponieważ jest on o~rząd wielkości krótszy niż proces hadronizacji ($10^{-24}$ sekundy), więc kwark $t$ nie tworzy żadnych złożonych cząstek. Po prostu rozpada się szybciej, niż może zostać związany w~postaci hadronów!

Jest to unikalna cecha kwarka wysokiego, który jako jedyna oddziałująca silnie cząstka elementarna może być obserwowany w~postaci swobodnej. Dzięki temu, że kwark $t$ nie zdąży wytworzyć hadronu, można bezpośrednio badać jego właściwości, np. spin. Obserwując kinematykę produktów rozpadu par kwark-antykwark ($t\bar{t}$), fizycy z~eksperymentów CMS i~ATLAS potrafili odtworzyć momenty pędu wyjściowych kwarków. Po analizie wielu takich zdarzeń znaleźli korelację między ich spinami -- co interpretujemy jako \Magenta{\textbf{splątanie kwantowe}} [1, 2].

Najnowsze odkrycia idą jednak jeszcze dalej. Choć kwark $t$ nie tworzy hadronów, to w~krótkiej chwili przed jego rozpadem możemy zaobserwować zaczątki hadronizacji. Oddziaływanie z~gluonami sprawia, że produkcja par $\bar{t}t$ jest zwiększona, co prowadzi do większej liczby obserwowanych zdarzeń, niż gdyby nie dochodziło do hadronizacji.

Takie zjawisko zaobserwowali niedawno fizycy pracujący w~eksperymencie CMS~[3]. Nową ,,cząstkę'' nazywamy \Magenta{\textbf{Toponium}} i~przedstawiamy jako nieuformowany stan związany kwarku wysokiego i~jego antycząstki $\bar{t}$. Jego masa jest równa około $340$ GeV, czyli jest niemal równa podwójnej masie kwarka $t$. Łatwo to zrozumieć -- wkład energetyczny od oddziaływań silnych (w~przeciwieństwie do lekkich hadronów, jak piony, protony czy neutrony) jest tu marginalny, ponieważ kwarki $t$ ($171$ GeV) są dużo cięższe niż suma energii zawartych w~Toponium gluonów (ok. $1$ GeV).
\marg[-30]{Warto zaznaczyć, że nie mamy jeszcze pełnego potwierdzenia, że nowo odkryta cząstka to na pewno Toponium. Alternatywy są jednak jeszcze bardziej ekscytujące -- oznaczałyby bowiem, że jest to zupełnie nowa cząstka elementarna, sygnalizująca odkrycie zupełnie nowej fizyki.\bigskip

  \textbf{Bibliografia}
  \vspace{4pt}
  
%\begin{dwieszpalty}
  \setbox0=\hbox{[0] }
  \hangindent\wd0
[1] ATLAS Collaboration.
\textit{Observation of quantum entanglement with top quarks at the ATLAS detector}.
Nature {633} (2024).
arXiv: 2311.07288 [hep-ex].

  \hangindent\wd0
[2] CMS Collaboration.
\textit{Measurements of polarization and spin correlation and observation of entanglement in top quark pairs using lepton + jets events from proton-proton collisions at $\sqrt{s} = 13$~TeV}.
Physical Review D {110} (2024).

  \hangindent\wd0
[3] CMS Collaboration.
\textit{Observation of a~pseudoscalar excess at the top quark pair production threshold}.
(eprint, marzec 2025)
arXiv: 2503.22382 [hep-ex].}

Ciekawym kierunkiem dalszych badań byłaby możliwość istnienia stanów wzbudzonych Toponium. W~przypadku lżejszych mezonów, takich jak~$J/\Psi$, znane są całe serie wyższych stanów energetycznych (np. $\Psi (2S)$), odpowiadających różnym modom oscylacyjnym czy orbitalnym kwarków. Podobne wzbudzenia, nawet w~postaci krótkich korelacji, moglibyśmy hipotetycznie zaobserować w~przypadku kwarka $t$.


Najważniejszy wydaje się jednak fakt, że w~LHC zarejestrowano zupełnie nową cząstkę. Jest to pierwsze tego typu odkrycie od zarejestrowania bozonu Higgsa, które miało miejsce ponad dziesięć lat temu. Istnienie Toponium stanowiło do tej pory hipotezę, i~wreszcie udało się ją potwierdzić eksperymentalnie. Jest to pierwszy hadron zawierający najcięższy z~kwarków, co czyni go wyjątkowym.

Nie ulega wątpliwości, że badania nad kwarkiem wysokim wciąż przynoszą nowe odkrycia. Możemy więc z~nadzieją i~ekscytacją czekać na przyszłość fizyki kwarka $t$.

%\advance\baselineskip by-.1pt%\vspace*{\parskip}
% \begin{dwieszpalty}
%\medskip


{}
\rightline{\large \it  Wiktor MATYSZKIEWICZ}
%\rightline{\large \it  Krzysztof TURZYŃSKI}
%\vspace*{-10pt}
%\vskip-10pt plus10pt
%\marg
{
}
%\end{dwieszpalty}
\eject
\normalsize